\Conclusion % заключение к отчёту
\label{cha:conclusion}

В ходе работы были решены следующие задачи:
\begin{enumerate}
    \item Разработан веб-сервер с модульной архитектурой и сквозным стеком вызовов от пользовательских HTML-форм до ядра СУБД PostgreSQL;
    \item Реализовано безопасное конструирование SQL-запросов с экранированием идентификаторов (\texttt{quoteIdent}) и литералов (\texttt{quoteLiteral});
    \item Реализована поддержка операций языка определения данных DDL: двухэтапное создание таблиц с поддержкой первичных и внешних ключей (включая самоссылающиеся ограничения), изменение структуры через ALTER TABLE, очистка через TRUNCATE и каскадное удаление через DROP TABLE CASCADE;
    \item Реализована поддержка операций языка манипулирования данными DML: пакетная вставка записей с автоматической обработкой NULL-значений, изменение данных по условию UPDATE и удаление данных DELETE;
    \item Реализован интерактивный конструктор запросов выборки SELECT с динамической проекцией, фильтрацией и сортировкой, а также оператор реляционного соединения INNER JOIN двух таблиц по общему ключу;
    \item Реализован механизм изоляции и быстрого восстановления исходного состояния базы данных через эндпоинт \texttt{POST /reset}.
\end{enumerate}

Корректность функционирования подтверждена серией тестов в браузере с прозрачным аудитом текста сформированных SQL-команд.
